\documentclass{article}
\usepackage{amsmath}
\usepackage{amssymb}
\usepackage{lmodern}

\title{Mathematical Physics: Calculus and Quantum Mechanics}
\author{RushTeX and XeLaTeX Verification Suite}
\date{September 2026}

\begin{document}

\maketitle

\section{Multivariable and Vector Calculus}

In classical field theory and differential geometry, integral theorems establish the duality between differential forms and boundary operators.

\subsection{Integral Theorems on Manifolds}
Green's Theorem in the plane relates a line integral around a simple closed curve $C$ to a double integral over the enclosed plane region $D$:
\[
\oint_C (L\,\mathrm{d}x + M\,\mathrm{d}y) = \iint_D \left( \frac{\partial M}{\partial x} - \frac{\partial L}{\partial y} \right) \mathrm{d}A
\]
Stokes' Theorem generalizes this to an oriented two-dimensional surface $\Sigma$ with boundary $\partial \Sigma$:
\[
\oint_{\partial \Sigma} \vec{A} \cdot \mathrm{d}\vec{r} = \iint_\Sigma (\nabla \times \vec{A}) \cdot \mathrm{d}\vec{S}
\]
The Gauss-Ostrogradsky Divergence Theorem over a compact three-dimensional region $V \subset \mathbb{R}^3$ with smooth boundary $\partial V$:
\[
\iiint_V (\nabla \cdot \vec{F})\,\mathrm{d}V = \iint_{\partial V} (\vec{F} \cdot \hat{n})\,\mathrm{d}S
\]

\subsection{Complex Analysis and Analytic Continuation}
The Cauchy Residue Theorem enables exact evaluation of meromorphic contour integrals along a positively oriented Jordan curve $\gamma$:
\[
\oint_\gamma f(z)\,\mathrm{d}z = 2\pi i \sum_{k=1}^n \mathrm{Res}(f, z_k)
\]
For Gaussian probability distributions and path integral normalizations, Euler's integral over the entire real line satisfies:
\[
\int_{-\infty}^\infty e^{-x^2}\,\mathrm{d}x = \sqrt{\pi}, \quad \int_{-\infty}^\infty e^{-\alpha x^2 + \beta x}\,\mathrm{d}x = \sqrt{\frac{\pi}{\alpha}} \exp\left( \frac{\beta^2}{4\alpha} \right)
\]

\subsection{Series Expansions and Asymptotics}
Taylor's Theorem with the Lagrange form of the remainder expresses analytic functions as:
\[
f(x) = \overbrace{\sum_{k=0}^n \frac{f^{(k)}(0)}{k!} x^k}^{\text{Taylor Polynomial}} + \underbrace{\frac{f^{(n+1)}(\xi)}{(n+1)!} x^{n+1}}_{\text{Lagrange Remainder}}
\]

\section{Quantum Physics and Operator Mechanics}

The state space of an isolated quantum system is represented by a complex separable Hilbert space $\mathcal{H}$ over $\mathbb{C}$.

\subsection{Schroedinger Wave Mechanics}
The dynamics of a non-relativistic quantum particle of mass $m$ subject to an external scalar potential $V(\vec{r}, t)$ obeys the Time-Dependent Schroedinger Equation:
\[
i \hbar \frac{\partial}{\partial t} \Psi(\vec{r}, t) = \left[ -\frac{\hbar^2}{2m} \nabla^2 + V(\vec{r}, t) \right] \Psi(\vec{r}, t)
\]
For stationary states $\Psi(\vec{r}, t) = \psi(\vec{r})e^{-iEt/\hbar}$, the Hamiltonian operator eigenvalue problem is:
\[
\hat{H} \psi_n = E_n \psi_n, \quad \hat{H} = -\frac{\hbar^2}{2m} \frac{\mathrm{d}^2}{\mathrm{d}x^2} + \frac{1}{2} m \omega^2 x^2
\]

\subsection{Heisenberg Uncertainty and Commutation Relations}
In the canonical operator formulation, the position operator $\hat{x}$ and momentum operator $\hat{p} = -i\hbar \frac{\partial}{\partial x}$ satisfy the Dirac commutator:
\[
[\hat{x}, \hat{p}] = \hat{x}\hat{p} - \hat{p}\hat{x} = i \hbar \mathbf{I}
\]
The Robertson-Schroedinger uncertainty relation for arbitrary observables $\hat{A}$ and $\hat{B}$ yields:
\[
\sigma_A \sigma_B \ge \frac{1}{2} \left| \langle [\hat{A}, \hat{B}] \rangle \right| \implies \sigma_x \sigma_p \ge \frac{\hbar}{2}
\]

\subsection{Dirac Bra-Ket Formalism and Spectral Decomposition}
In Dirac notation, transition amplitudes between state vectors $|\psi\rangle$ and $|\phi\rangle$ are given by the Hermitian inner product:
\[
\langle \psi | \phi \rangle = \int_{-\infty}^\infty \psi^*(x)\,\phi(x)\,\mathrm{d}x
\]
Resolution of identity in terms of complete orthonormal eigenstates:
\[
\sum_{n=0}^\infty |n\rangle \langle n| = \mathbf{I}, \quad \langle m | n \rangle = \delta_{mn}
\]

\subsection{Quantum Harmonic Oscillator and Ladder Operators}
Defining the dimensionless bosonic annihilation and creation operators $\hat{a}$ and $\hat{a}^\dagger$:
\[
\hat{a} = \sqrt{\frac{m\omega}{2\hbar}} \left( \hat{x} + \frac{i}{m\omega}\hat{p} \right), \quad \hat{a}^\dagger = \sqrt{\frac{m\omega}{2\hbar}} \left( \hat{x} - \frac{i}{m\omega}\hat{p} \right)
\]
Their fundamental commutator is $[\hat{a}, \hat{a}^\dagger] = 1$, allowing the Hamiltonian to be factorized into the number operator $\hat{N} = \hat{a}^\dagger \hat{a}$:
\[
\hat{H} = \hbar \omega \left( \hat{a}^\dagger \hat{a} + \frac{1}{2} \right), \quad E_n = \hbar \omega \left( n + \frac{1}{2} \right)
\]

\subsection{Spin Angular Momentum and Pauli Matrices}
For spin-$1/2$ systems (such as electrons), the observable spin vector $\vec{S} = \frac{\hbar}{2}\vec{\sigma}$ is represented by the hermitian $2 \times 2$ Pauli matrices:
\[
\sigma_x = \begin{pmatrix} 0 & 1 \\ 1 & 0 \end{pmatrix}, \quad \sigma_y = \begin{pmatrix} 0 & -i \\ i & 0 \end{pmatrix}, \quad \sigma_z = \begin{pmatrix} 1 & 0 \\ 0 & -1 \end{pmatrix}
\]
The mixed state statistical density operator $\rho$ satisfies the Liouville-von Neumann equation:
\[
\rho = \frac{1}{2}\left( \mathbf{I} + \vec{P} \cdot \vec{\sigma} \right), \quad i \hbar \frac{\partial \rho}{\partial t} = [\hat{H}, \rho]
\]

\subsection{Relativistic Quantum Mechanics}
The covariant Dirac equation for a spin-$1/2$ fermion in Minkowski space $(\mathbb{R}^{1,3}, \eta_{\mu\nu})$:
\[
\left( i \hbar \gamma^\mu \partial_\mu - mc \right) \psi = 0, \quad \{\gamma^\mu, \gamma^\nu\} = 2\eta^{\mu\nu}\mathbf{I}_{4\times 4}
\]

\end{document}
